\chapter{Uncertainty Quantification}
\label{ch:uncertainty}

This chapter defines the uncertainty machinery of \jns{}: how the committee
disagreement is measured, how the fallback threshold is calibrated, and how
the calibration quality is audited.  Everything here maps one-to-one onto
\code{jnuscout.calculators.committee} and \code{jnuscout.uq.calibration}.

\section{Disagreement measures}

The committee consists of ML potentials from different architecture
families (by default MACE-MP-0, MACE-MPA-0 and AIMNet2-b973c).  For a
geometry with $N$ atoms, let member $i$ predict forces $F_{i,j}$ on atom $j$
and energy $E_i$.

\subsection{Force disagreement}

The per-atom deviation of each member from the committee mean force is
mapped to its vector norm and maximised over all atoms and members:

\begin{equation}
  \mathrm{uq}_F
  = \max_{i,\,j}\;
    \left\lVert F_{i,j}
      - \frac{1}{M}\sum_{i'} F_{i',j} \right\rVert
  \qquad [\mathrm{eV}/\text{\AA}]
  \label{eq:uqf}
\end{equation}

The maximum, not the mean, is deliberate: in reacting geometries the
failure mode is typically a single atom with a wrong force direction, and a
mean over atoms would dilute it.

\subsection{Energy disagreement}

The standard deviation of the members' energies, divided by the atom count
to remove the size dependence:

\begin{equation}
  \mathrm{uq}_E = \frac{\operatorname{std}_i(E_i)}{N}
  \qquad [\mathrm{eV}/\mathrm{atom}]
  \label{eq:uqe}
\end{equation}

Energy disagreement is computed \emph{within one backend family} (by
default the largest family, the two MACE models).  Cross-family absolute
energies are not comparable -- MACE reports atomization energies, AIMNet2
reports absolute total energies -- so an optional per-member, per-atom
baseline offset (\code{energy_baselines}) can be supplied to align members
before a cross-family spread is even defined.  Without baselines, the raw
absolute-energy spread is diagnostics-only.  Degenerate inputs (keys that do
not match the members, a family filter without members, missing baselines)
emit \code{logging} warnings rather than errors: the numeric behaviour of
the tolerant fallback path is preserved, but silence is not.

\section{Threshold calibration}

A threshold $\theta$ converts the disagreement score into a decision.  The
calibration protocol implemented by \code{calibrate_threshold} is:

\begin{equation}
  \theta = \operatorname{Quantile}\!\left(\,
      \mathrm{uq} \;\middle|\; \mathrm{err} > \varepsilon_{\mathrm{crit}},
      \; 5\% \,\right)
  \label{eq:theta}
\end{equation}

\emph{Take the low quantile of the uncertainty score among the dangerous
configurations} -- those whose true error (against the reference
calculations) exceeds the tolerance $\varepsilon_{\mathrm{crit}}$.  Only
five percent of dangerous configurations sit below $\theta$, so the
false-negative rate of the fallback trigger is bounded by construction.

\begin{notebox}[Why not the 95th percentile of uq?]
A threshold taken from the score's own distribution describes the high tail
of the score -- which may fall between two modes of the score distribution
and catch dangerous configurations arbitrarily.  The conditional quantile in
Equation~\ref{eq:theta} asks the right question: \emph{among configurations
that are actually bad, how low can the score be?}
\end{notebox}

If the dangerous set is empty (no reference error exceeds
$\varepsilon_{\mathrm{crit}}$), the function returns \code{inf}, meaning
``never fall back'' -- the caller is expected to treat this as a signal that
the reference set was too easy, not as a validation of the potential.

\section{Coverage is a consequence, not a constant}

The fallback rate is not a fixed performance number of the method: it
follows from the accuracy requirement.  A stricter
$\varepsilon_{\mathrm{crit}}$ (more configurations labelled dangerous) forces
a lower $\theta$ and a higher fallback rate.  Quote the pair
$(\varepsilon_{\mathrm{crit}}, \text{fallback rate})$ together, never the
rate alone.

\section{Calibration API}

\begin{description}
  \item[\code{calibrate_threshold(uq, err, epsilon_crit, quantile=0.05)}]
    Returns $\theta$ per Equation~\ref{eq:theta}; raises \code{ValueError}
    on shape mismatch; returns \code{inf} for an empty dangerous set.
  \item[\code{spearman_rho(uq, err)}]
    Rank correlation between the score and the true error.  A usable
    committee should reach $\rho$ well above $0.6$; the value is reported,
    not thresholded, by the library.
  \item[\code{expected_calibration_error(uq, err, epsilon_crit, n_bins=10)}]
    Expected calibration error of the threshold decision over binned scores.
  \item[\code{coverage_stats(uq, err, theta, epsilon_crit)}]
    For a chosen threshold: number of configurations, fallback rate,
    number of dangerous configurations, false-negative rate, and the mean
    absolute error of the kept (non-fallback) predictions versus the full
    set.
\end{description}

\section{Measured record}

On the development systems of this project (C/H/N/O chemistry, ORCA L2 as
the reference), the committee disagreement reached a rank correlation of
$\rho = 0.847$ against true force errors at calibration, and the fallback
trigger fired preferentially in transition-state regions, where the
disagreement rose by several times over equilibrium values.  These numbers
characterise the protocol on the tested systems; they are not transferable
accuracy guarantees for other chemistries.

\begin{readbox}[Auditing a calibrated threshold]
A threshold is only as good as its calibration set.  Before trusting
$\theta$, check (i) that the reference set covers the force regimes the
exploration will visit (equilibrium \emph{and} stretched geometries),
(ii) the rank correlation $\rho$ of the score against the true error, and
(iii) the coverage statistics at the chosen
$\varepsilon_{\mathrm{crit}}$.  The utility script
\script{calibrate_uq} performs all three on a reference set it builds from
small molecules plus random displacements.
\end{readbox}
